\chapter{Background}\label{s:background}

This chapter provides the conceptual foundations necessary to understand the
remainder of the thesis. We cover scientific workflow systems, data provenance,
FAIR principles for computational workflows, intent-based orchestration, and the
use of LLMs in scientific computing. We conclude with an introduction to Brane
and BraneScript, the execution platform targeted by this work.

% ------------------------------------------------------------------
\section{Scientific Workflow Management Systems}
% ------------------------------------------------------------------

A \emph{scientific workflow} is a structured composition of computational steps
whose execution produces a reproducible scientific result. Workflow Management
Systems (WMSs) automate the orchestration of these steps across potentially
heterogeneous compute resources, handling scheduling, data transfer,
error recovery, and resource allocation transparently from the
scientist~\cite{deelman2015pegasus}.

\subsection{Representative Systems}

Several WMS platforms have achieved broad adoption in the scientific community.

\textbf{Pegasus}~\cite{deelman2015pegasus} is a DAG-based system developed at
USC/ISI that abstracts over heterogeneous grid and cloud resources. It supports
four distinct lifecycle phases: workflow generation, workflow planning
(site selection and data staging), workflow execution, and monitoring. Pegasus
uses HTCondor as its underlying execution engine and emphasises fault tolerance
through automatic retry and checkpointing.

\textbf{Nextflow}~\cite{di2017nextflow} adopts a dataflow programming model in
which independent computational steps (processes) communicate via asynchronous
channels. It has become dominant in bioinformatics due to its native Docker and
Singularity container support and its community registry of reusable pipelines
(nf-core). Nextflow's DSL2 syntax is Groovy-based and explicitly describes
data movement between processes.

\textbf{Snakemake}~\cite{molder2021snakemake} follows a \emph{make}-inspired
declarative model: scientists define rules that describe how to produce output
files from input files, and the system infers the necessary execution order.
It is particularly popular in genomics and integrates natively with conda
environments for software dependency management.

\textbf{Galaxy}~\cite{afgan2018galaxy} provides a web-based visual interface
for workflow construction, lowering the barrier for non-programmer scientists.
Workflows are composed graphically and can be shared through public repositories.

\textbf{Common Workflow Language (CWL)}~\cite{crusoe2022cwl} is a
vendor-neutral specification rather than an execution engine. It defines a
portable, YAML-based format for describing tool invocations and their data
connections, with the goal of enabling workflow portability across WMS
implementations.

Despite their diversity, all these systems share a common limitation from the
perspective of this thesis: they require the scientist to specify workflows
explicitly and at the computational level. The transition from a scientific
intent to an executable workflow is left entirely to the researcher.

\subsection{Workflow Lifecycle}

A typical WMS workflow follows five phases~\cite{deelman2015pegasus}:
(1)~\emph{specification}, where the scientist describes the computational steps
and their dependencies; (2)~\emph{planning}, where the WMS maps abstract steps
to concrete computational resources; (3)~\emph{execution}, where steps run on
target infrastructure; (4)~\emph{monitoring}, where execution progress and
failures are tracked; and (5)~\emph{archiving}, where results and provenance
are stored for future reference. This thesis focuses on the interface before
phase~1: the translation of natural language into a formal workflow specification.

% ------------------------------------------------------------------
\section{Data Provenance}
% ------------------------------------------------------------------

Reproducibility in scientific computing requires not only that results can be
recomputed, but that the precise computational history that produced those
results is recorded and queryable. This is the domain of \emph{data provenance}.

\subsection{W3C PROV Data Model}

The W3C PROV Data Model (PROV-DM) provides a standard vocabulary for expressing
provenance. Its core concepts are \emph{entities} (data artefacts),
\emph{activities} (computational steps), and \emph{agents} (processes or
people). Relationships such as \texttt{wasGeneratedBy}, \texttt{used}, and
\texttt{wasAttributedTo} connect these elements into a provenance graph that
captures the causal history of any output.

\subsection{Prospective vs.\ Retrospective Provenance}

Provenance is commonly distinguished along two axes. \emph{Prospective
provenance} captures the workflow specification---the plan of what will be
executed---and is equivalent to the workflow definition itself.
\emph{Retrospective provenance} captures what actually ran, including concrete
inputs, parameter values, timestamps, and intermediate results. Full
reproducibility requires both: the specification must be re-runnable
(prospective), and each execution must be uniquely identified (retrospective).

\subsection{Workflow Provenance Formats}

Several formats extend PROV-DM for workflows.
\textbf{CWLProv}~\cite{soiland2018cwlprov} records the full execution of a CWL
workflow in a Research Object (RO) bundle, capturing all inputs, intermediate
files, and PROV-O annotations. \textbf{RO-Crate}~\cite{soiland2022rocrate} is
a more lightweight packaging format based on schema.org that bundles data with
its metadata and provenance, enabling FAIR sharing of computational results.

Intent-based systems add a new provenance challenge: not only must the execution
be recorded, but the \emph{intent} that triggered the execution must also be
captured. This is what we term \emph{intent provenance}: the mapping from a
natural-language request to the workflow that was generated and run. None of the
existing provenance formats directly address this. This gap is one of the
motivating design considerations of the semantic cache described in this thesis.

% ------------------------------------------------------------------
\section{FAIR Principles for Workflows}
% ------------------------------------------------------------------

The FAIR principles---Findable, Accessible, Interoperable, and
Reusable---were originally articulated for scientific data, but have since been
extended to computational workflows~\cite{goble2020fair}. FAIR workflows require
that workflow specifications be discoverable in registries, accessible via
stable identifiers, portable across execution environments, and documented with
sufficient metadata to enable third-party reuse.

In practice, achieving FAIR compliance for LLM-generated workflows is
non-trivial: outputs are generated dynamically, may not be stored in a registry,
and may vary across sessions. The semantic cache introduced in this thesis
addresses the reusability dimension by ensuring that once a workflow is generated
for a given intent, subsequent requests for semantically equivalent intents
retrieve the same (or a semantically equivalent) workflow without additional
LLM inference.

% ------------------------------------------------------------------
\section{Intent-Based Orchestration}
% ------------------------------------------------------------------

\emph{Intent-based networking} (IBN) is a paradigm that originated in network
management, where operators specify \emph{what} a network should do rather than
\emph{how} to configure it~\cite{bn2016ibn}. The system then autonomously
translates the high-level intent into low-level configuration. The IBN lifecycle
comprises four phases: \emph{expression} (the operator states an intent),
\emph{translation} (the system converts the intent to a configuration),
\emph{activation} (the configuration is applied), and \emph{assurance} (the
system verifies that the intent is satisfied).

\subsection{Application to Scientific Computing}

The IBN paradigm generalises naturally to scientific computing. A scientist
expresses an intent (``compute the GC content of this sequence''), the system
translates it into a workflow, executes the workflow on the target
infrastructure, and verifies that the output is correct. This framing shifts the
cognitive burden from workflow engineering to intent expression, which is a
task scientists are already equipped to perform.

The MAPE-K (Monitor--Analyse--Plan--Execute with shared Knowledge) control loop
provides a complementary model for adaptive orchestration: the system monitors
execution, analyses failures, re-plans the workflow, and re-executes, using a
shared knowledge base to inform decisions. This model is relevant when LLM
outputs fail to compile or execute, as the system can iterate on the translation
rather than surfacing a raw error to the user.

\subsection{Gaps in the Literature}

A systematic review of intent-based orchestration for scientific computing
reveals five persistent gaps:

\begin{description}
    \item[G1 --- Intent Provenance Model.] No existing system defines a formal
        model for recording the intent that triggered a workflow execution
        alongside the standard retrospective execution provenance.

    \item[G2 --- Hybrid Intent-to-WMS Bridge.] No system simultaneously supports
        declarative natural-language intent expression, translation to a
        target WMS DSL, and WMS-grade reproducibility guarantees. Systems either
        provide natural-language interfaces without WMS integration, or WMS
        integration without natural-language interfaces.

    \item[G3 --- Intent Reproducibility Definition.] There is no agreed
        definition of what it means for an intent-based system to be
        reproducible: does the same intent need to produce the same workflow, or
        just a semantically equivalent one?

    \item[G4 --- Human-in-the-Loop Trust Checkpoints.] Existing proposals
        for LLM-based workflow generation do not provide principled mechanisms
        for a scientist to inspect and approve a generated workflow before
        execution, particularly important in safety-sensitive or privacy-sensitive
        domains.

    \item[G5 --- Cross-Paradigm Benchmarks.] There are no publicly available
        benchmarks that evaluate intent-based systems against conventional WMS
        approaches on the same scientific tasks.
\end{description}

This thesis addresses \textbf{G2} directly: it implements a hybrid bridge
between natural-language intent expression and the Brane WMS via LLM-based
BraneScript generation. It provides partial evidence towards \textbf{G3} through
the semantic cache threshold analysis, and motivates \textbf{G1} as future work.

% ------------------------------------------------------------------
\section{Large Language Models in Scientific Computing}
% ------------------------------------------------------------------

\subsection{LLMs for Code and Workflow Generation}

LLMs have demonstrated substantial capability in code generation tasks. Models
such as OpenAI Codex~\cite{chen2021evaluating} and its successors have been
evaluated on programming benchmarks including HumanEval and MBPP, where they
achieve competitive pass@k rates. More recent open-weight models (Qwen, Llama,
Mistral families) have closed much of the performance gap with proprietary
systems, making local deployment on HPC infrastructure feasible.

In the workflow domain, Balis~\cite{balis2024llm} demonstrated that LLMs can
generate Parsl-based Python workflows from natural-language descriptions,
achieving promising accuracy on domain-specific tasks. Shin et al.\
demonstrated similar results for bioinformatics pipelines. These works
established that LLMs can encode sufficient domain knowledge about workflow
patterns to serve as initial translators, but noted that accuracy degrades
significantly when the target WMS has a non-standard or less-common DSL---a
challenge directly relevant to BraneScript.

\subsection{Retrieval-Augmented Generation}

Retrieval-Augmented Generation (RAG)~\cite{lewis2020rag} augments LLM inference
by retrieving relevant context from an external knowledge base and prepending it
to the input prompt. In the workflow translation setting, the knowledge base
contains package documentation---function signatures, parameter descriptions, and
usage examples---and the retrieval system selects the subset of packages most
likely relevant to the current intent.

RAG addresses a critical limitation of pure LLM generation: models cannot know
the exact API of a custom execution environment such as Brane without having
seen it in training data. By retrieving package documentation at inference time,
the system grounds generation in the actual available capabilities, reducing
hallucinated function calls and incorrect parameter types.

\subsection{ControlA and Agentic Orchestration}

Gueroudji et al.~introduced ControlA, a system that uses LLMs as autonomous
agents for scientific workflow management, allowing the model to iteratively
refine a workflow in response to execution feedback. This agentic approach is
complementary to the single-shot translation studied in this thesis, and
represents a potential future extension in which the LLM can retry or repair
failed BraneScript workflows automatically.

\subsection{Limitations Relevant to This Work}

Two limitations of LLM-based workflow generation are especially relevant here.
First, sentence embedding models used for semantic similarity (such as the
\texttt{all-MiniLM-L6-v2} model used in the semantic cache) represent meaning
at the level of propositions, not at the level of specific parameter values.
Two intents that differ only in a numeric threshold or dataset name will embed
at very high similarity (often $>0.97$ cosine similarity), making it impossible
to distinguish them purely by similarity threshold. This is a fundamental
tension between semantic caching and scientific reproducibility that we analyse
empirically in Chapter~\ref{s:evaluation}.

Second, LLMs tend to hallucinate plausible but non-existent function calls when
the target DSL is underrepresented in training data. Fine-tuning on
domain-specific BraneScript examples is our primary mitigation, complemented by
RAG-based package retrieval.

% ------------------------------------------------------------------
\section{Brane and BraneScript}
% ------------------------------------------------------------------

\subsection{The Brane Framework}

Brane~\cite{brane2022} is a distributed, containerised computing framework
designed for federated scientific data analysis. It targets scenarios where
data cannot leave a site due to privacy or regulatory constraints
(e.g., hospital patient records, genomics data under GDPR), but compute results
can be returned to a central orchestrator. Brane wraps computational tools in
OCI-compatible container images (\emph{packages}), and provides a central
registry from which packages can be invoked across participating sites.

Each Brane package exposes a set of typed functions. A function call specifies
a package name, a function name, and typed argument values; Brane dispatches
the call to a site that holds the relevant data and returns the result to the
orchestrator. This model provides strong isolation and auditability: every
function invocation is logged, every data transfer is explicit, and every
result is versioned.

\subsection{BraneScript}

BraneScript is the DSL used to express Brane workflows. It is a statically
typed, procedural language with Python-inspired syntax. Key constructs include:

\begin{itemize}
    \item \textbf{Package imports}: \texttt{import \textit{package\_name};} loads
        a package's function signatures into scope.
    \item \textbf{Function calls}: \texttt{let result := \textit{func}(\textit{arg} := \textit{value});} invokes a package function with named arguments.
    \item \textbf{Data commits}: \texttt{commit\_result(\textit{name}, \textit{result});} persists an intermediate result as a named dataset.
    \item \textbf{Control flow}: Standard \texttt{if}/\texttt{else} and
        \texttt{for} constructs.
    \item \textbf{Parallel execution}: \texttt{parallel [\textit{stmt1}, \textit{stmt2}]} executes multiple statements concurrently.
\end{itemize}

The key challenge for LLM-based generation is that BraneScript is not
represented in standard pre-training corpora. Models must infer its syntax from
the few examples provided in the system prompt and from the RAG context.
Supervised fine-tuning on our curated dataset of 584 BraneScript examples is
intended to internalise the syntactic patterns and package invocation conventions
that the base model cannot reliably produce zero-shot.

\subsection{Design Principles}

Two design principles motivate the architecture of our system.
\emph{Compile-before-execute}: generated BraneScript is validated by the Brane
compiler before any execution is attempted, catching syntactic errors without
consuming compute resources. \emph{Progressive trust}: the system first attempts
cache retrieval (deterministic, zero-cost), then LLM generation (stochastic,
high-cost), and only executes after compilation succeeds. This ordering minimises
both latency and the risk of executing a malformed workflow.


